% Options for packages loaded elsewhere
% Options for packages loaded elsewhere
\PassOptionsToPackage{unicode}{hyperref}
\PassOptionsToPackage{hyphens}{url}
%
\documentclass[
  ignorenonframetext,
]{beamer}
\newif\ifbibliography
\usepackage{pgfpages}
\setbeamertemplate{caption}[numbered]
\setbeamertemplate{caption label separator}{: }
\setbeamercolor{caption name}{fg=normal text.fg}
\beamertemplatenavigationsymbolsempty
% remove section numbering
\setbeamertemplate{part page}{
  \centering
  \begin{beamercolorbox}[sep=16pt,center]{part title}
    \usebeamerfont{part title}\insertpart\par
  \end{beamercolorbox}
}
\setbeamertemplate{section page}{
  \centering
  \begin{beamercolorbox}[sep=12pt,center]{section title}
    \usebeamerfont{section title}\insertsection\par
  \end{beamercolorbox}
}
\setbeamertemplate{subsection page}{
  \centering
  \begin{beamercolorbox}[sep=8pt,center]{subsection title}
    \usebeamerfont{subsection title}\insertsubsection\par
  \end{beamercolorbox}
}
% Prevent slide breaks in the middle of a paragraph
\widowpenalties 1 10000
\raggedbottom
\AtBeginPart{
  \frame{\partpage}
}
\AtBeginSection{
  \ifbibliography
  \else
    \frame{\sectionpage}
  \fi
}
\AtBeginSubsection{
  \frame{\subsectionpage}
}
\usepackage{iftex}
\ifPDFTeX
  \usepackage[T1]{fontenc}
  \usepackage[utf8]{inputenc}
  \usepackage{textcomp} % provide euro and other symbols
\else % if luatex or xetex
  \usepackage{unicode-math} % this also loads fontspec
  \defaultfontfeatures{Scale=MatchLowercase}
  \defaultfontfeatures[\rmfamily]{Ligatures=TeX,Scale=1}
\fi
\usepackage{lmodern}

\usetheme[]{metropolis}
\ifPDFTeX\else
  % xetex/luatex font selection
\fi
% Use upquote if available, for straight quotes in verbatim environments
\IfFileExists{upquote.sty}{\usepackage{upquote}}{}
\IfFileExists{microtype.sty}{% use microtype if available
  \usepackage[]{microtype}
  \UseMicrotypeSet[protrusion]{basicmath} % disable protrusion for tt fonts
}{}
\makeatletter
\@ifundefined{KOMAClassName}{% if non-KOMA class
  \IfFileExists{parskip.sty}{%
    \usepackage{parskip}
  }{% else
    \setlength{\parindent}{0pt}
    \setlength{\parskip}{6pt plus 2pt minus 1pt}}
}{% if KOMA class
  \KOMAoptions{parskip=half}}
\makeatother


\usepackage{longtable,booktabs,array}
\usepackage{calc} % for calculating minipage widths
\usepackage{caption}
% Make caption package work with longtable
\makeatletter
\def\fnum@table{\tablename~\thetable}
\makeatother
\usepackage{graphicx}
\makeatletter
\newsavebox\pandoc@box
\newcommand*\pandocbounded[1]{% scales image to fit in text height/width
  \sbox\pandoc@box{#1}%
  \Gscale@div\@tempa{\textheight}{\dimexpr\ht\pandoc@box+\dp\pandoc@box\relax}%
  \Gscale@div\@tempb{\linewidth}{\wd\pandoc@box}%
  \ifdim\@tempb\p@<\@tempa\p@\let\@tempa\@tempb\fi% select the smaller of both
  \ifdim\@tempa\p@<\p@\scalebox{\@tempa}{\usebox\pandoc@box}%
  \else\usebox{\pandoc@box}%
  \fi%
}
% Set default figure placement to htbp
\def\fps@figure{htbp}
\makeatother





\setlength{\emergencystretch}{3em} % prevent overfull lines

\providecommand{\tightlist}{%
  \setlength{\itemsep}{0pt}\setlength{\parskip}{0pt}}



 


\setbeamerfont{title}{size=\large} \usepackage{hyperref} \setbeamertemplate{footline}[frame number]
\makeatletter
\@ifpackageloaded{caption}{}{\usepackage{caption}}
\AtBeginDocument{%
\ifdefined\contentsname
  \renewcommand*\contentsname{Table of contents}
\else
  \newcommand\contentsname{Table of contents}
\fi
\ifdefined\listfigurename
  \renewcommand*\listfigurename{List of Figures}
\else
  \newcommand\listfigurename{List of Figures}
\fi
\ifdefined\listtablename
  \renewcommand*\listtablename{List of Tables}
\else
  \newcommand\listtablename{List of Tables}
\fi
\ifdefined\figurename
  \renewcommand*\figurename{Figure}
\else
  \newcommand\figurename{Figure}
\fi
\ifdefined\tablename
  \renewcommand*\tablename{Table}
\else
  \newcommand\tablename{Table}
\fi
}
\@ifpackageloaded{float}{}{\usepackage{float}}
\floatstyle{ruled}
\@ifundefined{c@chapter}{\newfloat{codelisting}{h}{lop}}{\newfloat{codelisting}{h}{lop}[chapter]}
\floatname{codelisting}{Listing}
\newcommand*\listoflistings{\listof{codelisting}{List of Listings}}
\makeatother
\makeatletter
\makeatother
\makeatletter
\@ifpackageloaded{caption}{}{\usepackage{caption}}
\@ifpackageloaded{subcaption}{}{\usepackage{subcaption}}
\makeatother

\usepackage{bookmark}
\IfFileExists{xurl.sty}{\usepackage{xurl}}{} % add URL line breaks if available
\urlstyle{same}
\hypersetup{
  pdfauthor={Giorgio Arcara},
  hidelinks,
  pdfcreator={LaTeX via pandoc}}


\author{Giorgio Arcara}
\date{}

\begin{document}


\begin{frame}
% TITLE SLIDES
\title{Metodologia della Ricerca Clinica\\ \vspace{1em} \emph{09 - Scelta dei test per una valutazione}}
\author{Giorgio Arcara,\\ Università di Padova \\ IRCCS San Camillo, Venezia}

\titlegraphic{

\vspace*{7cm}
\includegraphics[scale=0.15]{Figures/LogoSanCamilloIRCCS_Unipd_alpha.png}
\hfill \includegraphics[scale=0.15]{Figures/CC_license_3_0.png}
}
\maketitle
\end{frame}

\begin{frame}{La scelta dei test}
\phantomsection\label{la-scelta-dei-test}
Due momenti di scelta dei test:

\begin{enumerate}
\item
  Scelta dei test da inserire nella propria ``cassetta degli attrezzi''
\item
  \textbf{Scelta dei test da utilizzare con uno specifico individuo per
  uno specifico scopo.}
\end{enumerate}
\end{frame}

\begin{frame}{La scelta dei test per una specifica valutazione}
\phantomsection\label{la-scelta-dei-test-per-una-specifica-valutazione}
Per la neuropsicologia clinica esistono tre principali approcci
(chiamati in molti modi in letteratura)

\begin{itemize}
\tightlist
\item
  Batteria fissa
\item
  Batteria flessibile
\item
  Batteria fissa di screening più test ad hoc
\end{itemize}
\end{frame}

\begin{frame}{Approccio a batteria fissa}
\phantomsection\label{approccio-a-batteria-fissa}
La batteria di test è la stessa per tutti gli individui valutati. In
genere è piuttosto ampia e tende a coprire tutti i domini cognitivi.

È molto diffusa specie in ambienti clinici in cui la valutazione tende
ad essere standardizzata per tutti i pazienti e in quei contesti in cui
gli individui valutati tendono ad essere abbastanza simili.
\end{frame}

\begin{frame}{Approccio a batteria fissa}
\phantomsection\label{approccio-a-batteria-fissa-1}
\textbf{PRO}

\begin{itemize}
\tightlist
\item
  È un approccio che non è influenzato da preconcetti del valutatore.
\item
  Permette una valutazione esaustiva del paziente.
\end{itemize}

\textbf{CONTRO}

\begin{itemize}
\tightlist
\item
  Richiede molto tempo (spesso più sedute) per completare una sola
  valutazione.
\item
  Può risultare molto faticoso per la persona che effettua i test.
\end{itemize}
\end{frame}

\begin{frame}{Approccio a batteria flessibile}
\phantomsection\label{approccio-a-batteria-flessibile}
I test vanno scelti di volta in volta sulla base di ipotesi
diagnositiche sull'esaminando. É accompagnata dal processo detto di test
delle ipotesi.

\emph{NOTA: in ambito forense, questo è l'approccio più diffuso, anche
per via dei tempi limitati della valutazione}
\end{frame}

\begin{frame}{Approccio a batteria flessibile}
\phantomsection\label{approccio-a-batteria-flessibile-1}
In genere si parta da un test, selezionato sulla base dei quesiti
diagnostici. In base al risultato del test si seleziona il test
successivo e così via.

Spesso è privilegiato un approccio che miri a falsificare (cioè
smentire) le proprie ipotesi, piuttosto che confermarli, secondo un
approccio deduttivo.

Es. se suppongo che il paziente abbia un deficit di memoria e ottengo
che nel test di memoria è nella norma ho falsificato la mia ipotesi e
posso dunque procedere con un'altra ipotesi, da indagare tramite un
altro test.
\end{frame}

\begin{frame}{Approccio a batteria flessibile}
\phantomsection\label{approccio-a-batteria-flessibile-2}
\textbf{PRO}

\begin{itemize}
\tightlist
\item
  È un approccio che permette di concentrarsi sulle specifiche richieste
  diagnostiche.
\item
  La durata è contenuta.
\end{itemize}

\textbf{CONTRO}

\begin{itemize}
\tightlist
\item
  L'esito può essere influenzato da preconcetti del neuropsicologo (es.
  bias di conferma).
\item
  Elementi cruciali dello stato cognitivo del paziente potrebbero essere
  trascurati.
\item
  Un approccio strettamente basato sulla falsificazione non tiene conto
  di possibili errori di misurazione.
\end{itemize}
\end{frame}

\begin{frame}{Approccio con batteria fissa di screening più test ad hoc}
\phantomsection\label{approccio-con-batteria-fissa-di-screening-piuxf9-test-ad-hoc}
È un approccio intermedio ai due. Si comincia con una valutazione
sommaria fatta con una batteria di screening a cui segue una valutazione
approfondita sulla base dei risultati allo screening.

\textbf{PRO} e \textbf{CONTRO} sono intermedi ai due approcci
precedenti.
\end{frame}

\begin{frame}{Altre considerazioni su scelta dei test}
\phantomsection\label{altre-considerazioni-su-scelta-dei-test}
Un'importante considerazione da fare sulla scelta dei test è che occorre
ricordarsi che (da un punto di vista statistico), più test si
somministrano, più è probabile che almeno uno vada al di sotto del
cut-off.

Una possibilità sarebbe quella di correggere i test con metodi di
correzione dei confronti multipli (es. Bonferroni). Nel metodo di
bonferroni per mantenere l'errore di 1° tipo al livello nominale (alpha
= 0.05, cioè 5\%) per tutti i test si aggiusto il valore critico
dividendo per il numero di test somministrati.

Supponiamo di avere somministrato 20 test, l'alpha critico diventerebbe
0.05/20 = 0.0025
\end{frame}

\begin{frame}{Altre considerazioni su scelta dei test}
\phantomsection\label{altre-considerazioni-su-scelta-dei-test-1}
Usare correzioni come il metodo di Bonferroni ha un altro problema:
assume che tutti i risultati dei test siano indipendenti. Sappiamo
invece che, venendo dallo stesso individuo, saranno verosimilmente
legati.

Esistono quindi metodi più corretti e che aumentano meno l'errore di 2°
tipo, di fatto non controllando completamente il 1°. Es. Benjamini \&
Yekutieli 2001, False Discovery Rate.

\emph{L'utilizzo di metodi di correzione di test è comunque molto poco
diffuso e potrebbe essere utilizzati solo con quei metodi che
restituiscono la probabilità esatta.}
\end{frame}




\end{document}
